% ============================================================================
\section{Método y alcance}\label{sec:metodo}

\subsection{De las personas a los escenarios}\label{sec:plantilla}

Un escenario es, siguiendo el material de apoyo del curso, una historia sobre una de las
personas ya creadas, que describe cómo esa persona completa una tarea o se comporta en una
situación determinada; proporciona un contexto, tiene actores, objetivos y una secuencia de
eventos, y cierra con un resultado (Baxter, Courage y Caine, 2015). Los mismos autores
distinguen con claridad los tres documentos que el equipo ha producido hasta ahora, y esa
distinción explica por qué esta actividad no repite a la anterior.

\begin{uxtabla}{L{2.6cm}Z{1.0}Z{1.0}}{Diferencia entre perfil de usuario, persona y escenario según el material de apoyo}
\uxheadrow \thd{Documento} & \thd{Qué es} & \thd{Para qué sirve} \\ \midrule
Perfil de usuario & Descripción detallada de los atributos de los usuarios & Saber para quién se construye el producto y a quién reclutar para investigar \\
Persona & Individuo ficticio creado para describir al usuario típico a partir del perfil & Mantener el foco en usuarios específicos durante las discusiones de diseño \\
Escenario & Historia que describe cómo una persona concreta completa una tarea en una situación dada & Poner a prueba el sistema, verificar que satisface las necesidades reales y construir funcionalidad que los usuarios efectivamente quieran usar \\
\end{uxtabla}

Los seis escenarios de este documento siguen una plantilla única, propuesta por McInerney
(2003) y recogida en el capítulo de referencia del curso. La plantilla ordena la historia en
cuatro secciones: un título que describe la situación y no al personaje, la situación y la
tarea, el método con el que el usuario aborda la situación, redactado de forma genérica y
neutral respecto de la tecnología, y la ruta de ejecución, donde ya se nombran las
funcionalidades concretas del producto. La tabla siguiente muestra cómo esa plantilla
contiene los cuatro elementos que solicita la actividad.

\begin{uxtabla}{L{3.4cm}Z{1.0}Z{1.0}}{Correspondencia entre la plantilla de escenario y los elementos solicitados}
\uxheadrow \thd{Sección de la plantilla} & \thd{Contenido} & \thd{Elemento solicitado que satisface} \\ \midrule
Contexto y usuario & Quién es, a qué audiencia pertenece, qué necesita, qué persigue y qué le duele en esta tarea & Definición del usuario \\
Expectativas del usuario & Qué espera encontrar, comprender o controlar al utilizar el producto en esta situación & Expectativas generadas \\
Situación y tarea & Situación inicial, reto y meta, sin anticipar todavía cómo se resuelve & Descripción de la tarea: qué quiere conseguir \\
Método para abordar la situación & Secuencia de pasos genérica, sin suponer un diseño concreto & Descripción de la tarea: pasos para lograrlo \\
Ruta de ejecución & Narrativa de cómo se completa la tarea en el portal, nombrando cada punto de contacto & Identificación de las interacciones clave \\
Resultado y sensación & Qué obtiene el usuario y cómo se siente al conseguirlo & Descripción del resultado deseado \\
Excepciones & Situaciones poco frecuentes pero relevantes, y cómo responde el producto & Elemento adicional recomendado por el material de apoyo \\
\end{uxtabla}

El método se mantiene neutral respecto de la interfaz, mientras que la ruta de ejecución
nombra los puntos de contacto y podrá actualizarse con el prototipo. Además, cada escenario
incluye excepciones relevantes, como recomienda el material de apoyo, para observar cómo
responde el producto fuera del flujo esperado.

\subsection{Cobertura de los escenarios}

Cada integrante desarrolló los escenarios de sus personas de la actividad anterior. En
conjunto cubren seis flujos principales y varían en contexto, urgencia, dispositivo y modo de
fallo.

\begin{uxtabla}{C{0.9cm}L{2.9cm}Z{1.0}Z{1.0}}{Cobertura de los seis escenarios sobre las personas y los flujos del producto}
\uxheadrow \thd{\#} & \thd{Persona} & \thd{Situación} & \thd{Flujo del producto que ejercita} \\ \midrule
1 & Laura Méndez & Validación urgente antes de una reunión & Acceso, buscador unificado, catálogo y asistente \\
2 & Diego Hernández & Cruce de tres fuentes y extracción reproducible & Explorador analítico y exportación en segundo plano \\
3 & Arturo Castañeda & Entender una señal de riesgo antes de exponer & Tableros consolidados, profundización progresiva y resumen asistido \\
4 & Roberto Valdez & Publicar una nueva versión del catálogo sin romper a nadie & Gobierno de metadatos, versiones y notificación dirigida \\
5 & Mariana Ovalle & Otorgar el acceso mínimo suficiente y poder retirarlo & Administración de usuarios, roles y evidencia \\
6 & Ximena Solís & Integrar un servicio y sobrevivir a un cambio de esquema & Consumo por API, credenciales de autoservicio y contrato de datos \\
\end{uxtabla}

Elena Ruiz, del perfil de riesgos y auditoría, y Jorge Mendieta, del perfil de ingeniería de
datos, se retomarán en la siguiente actividad, centrada en la arquitectura de información.

\subsection{Del escenario al journey map}

Las instrucciones piden descomponer el objetivo de un escenario en actividades y convertirlo
en un journey map. La tabla resume los mapas producidos y distingue su autoría.

\begin{uxtabla}{L{3.0cm}L{3.1cm}Z{1.0}Z{1.0}}{Los cuatro journey maps del entregable, su autoría y su objetivo}
\uxheadrow \thd{Mapa} & \thd{Elaboración} & \thd{Persona y objetivo} & \thd{Por qué ese recorrido} \\ \midrule
Individual & Jacqueline Sarmiento & Diego Hernández: cruzar tres fuentes y dejar el análisis reejecutable & Recorrido largo, con conciliación de datos y espera de proceso \\
Individual & Alexandro Mayoral & Arturo Castañeda: entender una señal de riesgo antes de exponer & Recorrido brevísimo y de máxima presión, en tableta y en quince minutos \\
Individual & Arthur Zizumbo & Ximena Solís: automatizar por API y sobrevivir a un cambio de esquema & Recorrido de semanas, donde lo decisivo ocurre después del uso \\
\textbf{De equipo} & \textbf{Los tres, en taller} & \textbf{Laura Méndez: validar una cifra y responder con la fuente citada} & \textbf{Persona primaria y denominador común de los demás perfiles} \\
\end{uxtabla}

Cada mapa cubre un objetivo acotado, como recomienda Walter (2022). Los recorridos
individuales contrastan duración, dispositivo y tipo de fricción; el
apartado~\ref{sec:entrelazados} muestra sus puntos de encuentro.

\subsection{Alcance declarado y plan de validación}\label{sec:alcance}

\begin{uxnota}[Alcance del método]
Esta primera versión es un mapa de supuestos, construido con la caracterización de la
audiencia, la experiencia profesional del equipo y los artefactos anteriores, todavía sin
datos de campo. Se utilizará como punto de partida para la investigación, no como evidencia
definitiva de la experiencia real.
\end{uxnota}

Los supuestos más fuertes del mapa se concentran en tres etapas y son los primeros que se
contrastarán con los instrumentos diseñados en la actividad anterior, cuya aplicación está en
curso:

\begin{uxlist}
\lead{Etapa 3, búsqueda}{se supone que el usuario formula su necesidad en lenguaje de negocio
  y no con el nombre técnico del campo. Lo verifica la pregunta abierta del instrumento sobre
  a quién o a qué recurre primero cuando no sabe dónde vive un dato.}
\lead{Etapa 4, validación}{se supone que la dependencia de un tercero para confirmar la
  definición de una variable es el principal consumidor de tiempo de la tarea. Lo verifica la
  pregunta sobre horas semanales dedicadas a buscar o validar un dato.}
\lead{Etapa 7, reutilización}{se supone que los usuarios no se enteran de que una definición
  cambió hasta que dos áreas presentan cifras distintas. Es el supuesto con menos evidencia
  documental y el que la entrevista semiestructurada debe explorar con mayor cuidado.}
\end{uxlist}

Las cifras de tiempo que aparecen en los escenarios están rotuladas como estimaciones del
equipo a partir de su experiencia profesional, o bien marcadas como pendientes de medición y
enlazadas a la pregunta del instrumento que las levantará. Ninguna cifra de este documento
procede de una fuente que el equipo no haya consultado directamente.

